\chapter{ARSITEKTUR SISTEM}

\section{Diagram Arsitektur}
Gambar~\ref{fig:architecture} memperlihatkan arsitektur sistem dengan pengelompokan tanggung jawab, bentuk storage yang berbeda, serta dependency antarservice. Ini adalah diagram arsitektur logis, bukan notasi UML murni atau diagram deployment jaringan. Panah putus-putus menunjuk dari pemakai layanan ke penyedianya; jalur event dijelaskan pada Bab 8.

% Layout inspired by https://tex.stackexchange.com/a/457624;
% shapes/layout authored for the actual IF4031 implementation.
\usetikzlibrary{fit,positioning,shapes.geometric,shapes.multipart}
\definecolor{archink}{HTML}{31465B}
\definecolor{archlight}{HTML}{F0F3F6}
\definecolor{archwarm}{HTML}{FBF4E6}
\begin{figure}[h]
\centering
\resizebox{\textwidth}{!}{%
\begin{tikzpicture}[
  app/.style={draw=archink,rounded corners=2pt,fill=white,align=center,text width=3.8cm,minimum height=1.6cm,font=\fontsize{9}{11}\selectfont,inner sep=5pt},
  store/.style={cylinder,shape border rotate=90,aspect=0.25,draw=archink,fill=archlight,align=center,text width=3.5cm,minimum height=1.9cm,font=\fontsize{9}{11}\selectfont,inner sep=7pt},
  bus/.style={app,rectangle,rounded corners=0pt,fill=archwarm,line width=0.65pt},
  actor/.style={app,fill=archlight,text width=12.5cm,minimum height=1.05cm},
  dep/.style={dashed,-{Straight Barb[length=2mm,width=1.6mm]},draw=archink,line width=0.65pt},
  lab/.style={font=\scriptsize,fill=white,inner sep=1.5pt},
  group/.style={draw=archink!45,dashed,rounded corners=4pt,inner sep=8pt},
  heading/.style={font=\scriptsize\bfseries,text=archink,fill=white,inner sep=2pt}
]
\node[actor] (client) at (0,6.8) {\textbf{Client downstream}\\Media/Pers \quad Tim Lapangan \quad BNPB Pusat/Internal};
\node[app] (api) at (0,4.2) {\textbf{Client-Facing API}\\Go :8080\\JWT / scope};
\node[app] (auth) at (5.3,4.2) {\textbf{Auth / Token Service}\\Go :8084\\RS256 / SQLite privat};
\node[app] (agg) at (0,0.2) {\textbf{BNPB Aggregator}\\Go :8083\\Ingest / query / outbox};
\node[app] (bmkg) at (-5.5,0.9) {\textbf{Mock BMKG}\\Go :8081 / HTTP\\Gempa / tsunami};
\node[app] (pvmbg) at (-5.5,-1.4) {\textbf{Mock PVMBG}\\Go :8082 / HTTP\\Volkanik v1 / v2};
\node[store] (pg) at (5.3,0.2) {\textbf{PostgreSQL}\\Canonical + history\\JSONB / outbox};
\node[bus] (nats) at (0,-3.8) {\textbf{NATS JetStream}\\Stream \texttt{HAZARD\_EVENTS}\\Persistent event broker};
\node[app] (dash) at (-5.3,-6.2) {\textbf{Dashboard Consumer}\\Durable / SQLite privat};
\node[app] (notify) at (0,-6.2) {\textbf{Notifier Consumer}\\Durable / SQLite privat};
\node[app,dashed] (third) at (5.3,-6.2) {\textbf{Third Consumer}\\Opsional: profil demo\\Durable / SQLite privat};
\node[group,fit=(api)(auth),label={[heading]above:Akses dan autentikasi}] {};
\node[group,fit=(bmkg)(pvmbg),label={[heading]above:Sumber simulasi}] {};
\node[group,fit=(agg)(pg)] (ingestgroup) {};
\node[heading,anchor=south east] at (ingestgroup.north east) {Ingest dan penyimpanan kanonik};
\node[group,fit=(nats)(dash)(notify)(third),label={[heading]below:Distribusi event dan pemrosesan independen}] {};
\draw[dep] (client.south) -- node[lab,right,pos=0.18]{REST / Bearer JWT} (api.north);
\draw[dep] (client.south -| auth.north) -- node[lab,right,pos=0.18]{login / refresh} (auth.north);
\draw[dep] (api.south) -- node[lab,left,pos=0.28]{HTTP / internal key} (agg.north);
\draw[dep] (agg.west) -- (-2.75,0.2) |- (bmkg.east);
\draw[dep] (agg.west) -- (-2.75,0.2) |- (pvmbg.east);
\draw[dep] (agg.east) -- node[lab,above]{SQL} (pg.west);
\draw[dep] (agg.south) -- node[lab,right,pos=0.5]{publish / PubAck} (nats.north);
\draw[dep] (dash.north) |- node[lab,above,pos=0.65]{pull / ACK} ([xshift=-1.1cm,yshift=-0.65cm]nats.south) -- ([xshift=-1.1cm]nats.south);
\draw[dep] (notify.north) -- (nats.south);
\draw[dep] (third.north) |- node[lab,above,pos=0.65]{pull / ACK} ([xshift=1.1cm,yshift=-0.65cm]nats.south) -- ([xshift=1.1cm]nats.south);
\draw[draw=archink,line width=0.5pt] ([xshift=-2pt,yshift=-2pt]nats.south west) rectangle ([xshift=2pt,yshift=2pt]nats.north east);
\end{tikzpicture}%
}
\caption{Arsitektur logis milestone 1. Kotak menunjukkan service; silinder adalah Canonical Store dan kotak ganda adalah broker. Third Consumer bersifat opsional. Bingkai mengelompokkan tanggung jawab, bukan batas Docker network. Panah consumer menuju broker menunjukkan dependency layanan, bukan arah delivery event.}
\label{fig:architecture}
\end{figure}

\section{Komponen dan Tanggung Jawab}
\begin{itemize}
  \item \textbf{Mock BMKG/PVMBG}: service pasif yang hanya merespons polling. Endpoint data mengikuti kontrak tugas; endpoint admin tambahan mendukung simulasi delay dan outage.
  \item \textbf{Aggregator}: menjalankan poller BMKG dan PVMBG secara terisolasi, memvalidasi payload, menggabungkan tsunami warning dengan event gempa, memetakan data, menyimpan canonical record, serta mengekspos API internal untuk query Client-Facing API.
  \item \textbf{Client-Facing API}: memeriksa scope, meminta data melalui API Aggregator, lalu menyaring field di server. Komponen ini tidak membaca PostgreSQL secara langsung.
  \item \textbf{Auth/Token Service}: menjaga identitas client yang berbeda, JWT RS256 berumur pendek dan state refresh pada SQLite privat. Client API memvalidasi JWT secara lokal dengan public key; permintaan baca tidak membutuhkan panggilan Auth.
  \item \textbf{NATS JetStream dan consumer}: Aggregator menerbitkan event; consumer dashboard dan notifikasi memiliki durable consumer masing-masing pada stream yang sama.
\end{itemize}
Pemisahan ini menjaga batas kepemilikan storage dan memungkinkan satu container dihentikan atau dibangun ulang tanpa menggabungkan seluruh sistem menjadi satu unit deploy [1].

\section{Alasan Pembagian Service}
Aggregator menyatukan polling, mapping, dan relay karena semuanya menggunakan Canonical Store yang sama. Loop dan pool koneksi masing-masing pekerjaan dipisah, tetapi seluruhnya tetap berhenti jika container Aggregator mati. Memisahkan ingest per instansi dapat memberi isolasi deployment tambahan, dengan biaya antarmuka dan koordinasi baru.

Client API dipisahkan dari Auth agar proses pembacaan cukup memegang public key, sedangkan private key penandatanganan tetap pada Auth. Gangguan Auth tidak membatalkan access token yang masih valid karena signature diperiksa secara lokal; login dan refresh tetap membutuhkan Auth. Pemisahan ini menambah panggilan HTTP dari API ke Aggregator, tetapi menjaga ownership database dan memisahkan perubahan akses client dari ingest.

Dashboard dan Notifier mempunyai container, durable, serta storage terpisah. Jika keduanya memakai satu durable, pesan dibagi sebagai pekerjaan, bukan disalin untuk setiap aplikasi. Durable terpisah memenuhi fan-out dan memungkinkan masing-masing mengejar backlog. Third Consumer memakai binary Notifier dengan durable dan volume sendiri, melalui profil \texttt{demo}; producer tidak perlu mengetahui consumer tambahan.

\clearpage
\section{Model Kelas Data}
Gambar~\ref{fig:domain-model} menggunakan UML Class Diagram untuk kontrak JSON, bukan struktur tabel database. Association gempa--warning mengikuti identifier sumber. Dependency \guillemotleft map\guillemotright\ dan \guillemotleft enrich\guillemotright\ adalah label pemetaan lokal, bukan inheritance. \texttt{DateTime} diserialisasikan sebagai string timestamp; \texttt{JSONValue} dapat berupa object, array, angka, string, boolean, atau null.

\usetikzlibrary{shapes.multipart}
\begin{figure}[h]
\centering
\resizebox{\textwidth}{!}{%
\begin{tikzpicture}[
  cls/.style={rectangle split,rectangle split parts=2,rectangle split part fill={archlight,white},draw=archink,align=left,text width=4.35cm,font=\fontsize{8}{10}\selectfont,inner sep=5pt,anchor=north},
  mapping/.style={dashed,-{Straight Barb[length=1.8mm,width=1.5mm]},draw=archink,line width=0.6pt},
  association/.style={draw=archink,line width=0.6pt},
  lab/.style={font=\fontsize{8}{9}\selectfont,fill=white,inner sep=1.5pt}
]
\node[cls] (seismic) at (-5.3,0) {\textbf{SeismicEvent}\nodepart{second}\ttfamily
 event\_id: String\\
 magnitude: Real\\
 depth\_km: Real\\
 epicenter\_lat: Real\\
 epicenter\_lon: Real\\
 region\_name: String\\
 occurred\_at: DateTime\\
 potential\_tsunami: Boolean};
\node[cls] (warning) at (0,0) {\textbf{TsunamiWarning}\nodepart{second}\ttfamily
 warning\_id: String\\
 related\_event\_id: String\\
 threat\_level: String\\
 affected\_zones: String [0..*]\\
 estimated\_arrival: DateTime};
\node[cls] (volcano) at (5.3,0) {\textbf{VolcanicReport}\nodepart{second}\ttfamily
 report\_id: String\\
 volcano\_id: String\\
 alert\_level: String\\
 eruption\_count\_24h: Integer\\
 ash\_column\_height\_m: Real\\
 reported\_at: DateTime\\
 confidence\_level: Real [0..1]};
\node[cls,text width=6.1cm] (hazard) at (-3.0,-5.1) {\textbf{HazardEvent}\nodepart{second}\ttfamily
 hazard\_id: UUID\\
 source: Source\\
 source\_ref\_id: String\\
 hazard\_type: HazardType\\
 severity: Severity\\
 area\_name: String\\
 latitude: Real\\
 longitude: Real\\
 occurred\_at: DateTime\\
 ingested\_at: DateTime\\
 attributes: Map<String, JSONValue>};
\node[cls,text width=4.85cm] (envelope) at (4.0,-5.1) {\textbf{Envelope}\nodepart{second}\ttfamily
 schema\_version: Integer\\
 event\_type: String\\
 message\_id: String\\
 correlation\_id: String\\
 published\_at: DateTime\\
 producer: String\\
 revision: Integer\\
 data: HazardEvent};
\draw[association] (seismic.north) -- ++(0,0.75) -| (warning.north);
\node[lab,left] at (-5.3,0.22) {1};
\node[lab,right] at (0,0.22) {0..*};
\node[lab,above] at (-2.65,0.75) {related\_event\_id};
\draw[mapping] (seismic.south) -- node[lab,left]{\guillemotleft map\guillemotright} (hazard.north -| seismic.south);
\draw[mapping] (warning.south) -- (0,-4.0) -- node[lab,above]{\guillemotleft enrich\guillemotright} (-3,-4.0) -- (hazard.north);
\draw[mapping] (volcano.south) -- (5.3,-4.65) -- node[lab,above]{\guillemotleft map\guillemotright} (-0.7,-4.65) -- (hazard.north -| -0.7,0);
\draw[association,-{Straight Barb[length=1.8mm,width=1.5mm]}] ([yshift=-1.7cm]envelope.north west) -- node[lab,above]{data} node[lab,below,pos=0.8]{1} ([yshift=-1.7cm]hazard.north east);
\end{tikzpicture}%
}
\caption{Kontrak sumber, entity kanonik, dan envelope event. Satu warning merujuk satu gempa; satu gempa dapat mempunyai banyak warning. \texttt{confidence\_level} bersifat opsional dan muncul pada PVMBG v2. \texttt{Envelope.data} berisi satu snapshot HazardEvent, bukan referensi ke latest row.}
\label{fig:domain-model}
\end{figure}

\noindent Enum kanonik: \texttt{Source} = BMKG/PVMBG; \texttt{HazardType} = SEISMIC/VOLCANIC; \texttt{Severity} = NORMAL/WASPADA/SIAGA/AWAS. Field tambahan sumber dipertahankan dalam \texttt{attributes}; metadata revisi dan delivery tetap berada di envelope. Warning yang belum memiliki gempa terkait dapat berada dalam snapshot sumber, tetapi tidak menjadi hazard mandiri.

\section{Protokol dan Alur Komunikasi}
\begin{enumerate}
\item Poller BMKG mengambil \texttt{/seismic-events} dan \texttt{/tsunami-warnings} secara concurrent. Poller PVMBG mengambil \texttt{/volcanic-reports} secara independen. Keduanya menggunakan full snapshot, bukan cursor \texttt{since}.
\item Aggregator memvalidasi dan memetakan data, lalu menyimpan latest hazard, revisi, dan outbox dalam satu transaksi. Relay menandai pesan terpublikasi setelah PubAck JetStream.
\item Client API memvalidasi JWT secara lokal dan membaca API internal Aggregator. Media hanya menerima ringkasan tanpa koordinat presisi, identifier sumber, atau attributes.
\item Consumer membaca stream secara independen dan ACK setelah pemrosesan berhasil. Detail mapping dan delivery ada pada Bab 4 dan 8.
\end{enumerate}

\section{Kepemilikan Storage dan Container}
PostgreSQL adalah Canonical Store milik Aggregator. Data terbaru berada pada \texttt{hazard\_events}; riwayatnya berada pada \texttt{hazard\_versions}. Tabel \texttt{source\_state} menyimpan freshness, dan \texttt{outbox} menyimpan envelope delivery. Field kanonik bertipe; attributes memakai JSONB. Auth memiliki SQLite untuk state refresh; setiap consumer memiliki SQLite untuk inbox dan side effect. Service lain tidak membuka koneksi ke Canonical Store. Sembilan service utama berjalan dalam container terpisah melalui Compose.

\section{Alur Permintaan dan Kegagalan}
Polling berjalan di background, bukan saat request client. Delay atau outage PVMBG tidak menahan BMKG maupun API baca. Data terakhir dapat ditandai stale/unavailable; full snapshot setelah pemulihan memperbarui freshness tanpa restart. Kegagalan broker menyisakan pending outbox; kegagalan consumer tidak menahan publish. Durable consumer dapat menerima redelivery selama pesan masih dalam retensi.

\subsection{Batas Jaringan}
Compose memisahkan \texttt{source-net}, \texttt{api-net}, \texttt{store-net}, dan \texttt{event-net}. Dua jaringan terakhir bersifat internal dan PostgreSQL/NATS tidak memublikasikan port host. Port aplikasi demo terikat ke \texttt{127.0.0.1}; deployment lintas host memerlukan TLS.
